%auto-ignore
\documentclass[main.tex]{subfiles}

\begin{document}

\section{The Free Independence Principle of \netsort{} Programs}
\label{sec:FIPMain}

A powerful analogue of \emph{independence} (of scalar random variables) in random matrix theory is called \emph{Asymptotic Free Independence}, or \emph{Asymptotic Freeness}, defined below in \cref{defn:asfree}.
\begin{defn}\label{defn:asfree}
    Fix $k$.
Consider collections of random matrices $\mathcal{W}_{n}^{1},\ldots,\mathcal{W}_{n}^{k}\subseteq\R^{n\times n}$ for each $n\ge1$, of constant cardinalities (with $n$).
(For example, each $\mathcal{W}_{n}^{i}$ can be $\{W, W^{\trsp}\}$ for some weight matrix $W$ in a neural network).
We say $\mathcal{W}_{n}^{1},\ldots,\mathcal{W}_{n}^{k}$ are \emph{almost surely asymptotically free}%
\footnote{
    We can also speak of asymptotically free in expectation, in which case, we want the expectation of the trace to converge to 0.
    In most scenarios (where we have tail bounds on the matrix operator norms), this is a weaker notion than almost sure asymptotic freeness.
}, if
\[
\inv n \tr\left(\prod_{i=1}^{k}\left(P_{i}(\mathcal{W}_{n}^{j_{i}})-\tau_{i}\right)\right)\asto 0
\]
where $\tau_{i}=\inv n \tr(P_{i}(\mathcal{W}_{n}^{j_{i}}))$, $P_{i}$ is a (noncommutative) polynomial in $|\mathcal W^{j_i}_n|$ variables, and $j_{1},\ldots,j_{k}\in[k]$ are indices with no two adjacent $j_{i}$ equal,
with $\{P_i\}_i$, $\{j_i\}_i$ independent of $n$.%
\footnote{
    Note here $\prod$ is a non-commutative product, where in its expansion, $i$ goes from right to left.
}
\footnote{One can also formulate the asymptotic freeness of subalgebras of a non-commutative algebra, in which case our definition here is equivalent to the freeness of the subalgebras generated by the respective collections of random matrices.}
\end{defn}

Whereas the independence of \emph{scalar} random variables allows one to compute the expectation of their sum and product easily, the asymptotic freeness of random \emph{matrices} allows one to compute the asymptotic spectral distributions of their sum and (matrix) product easily.
Independence implies that two random variables are in ``general position'' and thus cannot conspire to fluctuate in the same direction. Likewise, asymptotic freeness of two random matrices, intuitively, implies that their respective eigenvectors and eigenvalues lie in general position to each other, so that their spectral distributions would combine in predictable ways were they to be summed or multiplied. For more background, see \citep{tao_topics_2012}.

\emph{A priori}, because the activations and preactivations of a neural network depend in a highly nonlinear and complex manner in the weight matrices (and biases), one can hardly suspect that activations can be ``independent'' from the weight matrices in some way. 
However, we will show, in a randomly initialized neural network of any architecture, the weight matrices are asymptotically free from (the diagonal matrices formed from) the preactivations of the network.
This will follow from the much more general \emph{Free Independence Principle} of \netsort{} programs, \cref{thm:FreeIndependencePrinciple}.
We give a proof in \cref{{sec:FIP}} that follows the same overall strategy as the proof of the Semicircle Law in \cref{sec:semicircle}.
\begin{thm}[Free Independence Principle, for Tensor Programs]
\label{thm:FreeIndependencePrinciple} Consider any \netsort{} program $\pi$ with polynomially bounded nonlinearities
and \cref{setup:netsort}. Then the random matrix collections $\{W,W^{\trsp}\}$ for every $W\in\mathcal{W}$, along with the collection of diagonal matrices $\mathcal D(\pi)$ (defined immediately below in \cref{defn:DMatrices}) are asymptotically free as $n\to\infty$.
\end{thm}
\begin{defn}\label{defn:DMatrices}
    Suppose $\mathcal X$ is a subset of ($\R^n$) vectors in a program.
    Let $\mathcal D(\mathcal X)$ denote the (infinite) collection of diagonal matrices formed from bounded, coordinatewise images of $\mathcal X$:
    \begin{equation}\mathcal D(\mathcal X) \defeq \{\Diag(\psi(x^1,\ldots, x^k)): k \ge 0;\ \ x^1, \ldots, x^k \in \mathcal X;\ \ \psi: \R^k \to \R \text{ bounded}\} \sbe \R^{n\times n}.
        \label{eqn:DMat}
    \end{equation}
    If $\pi$ is a program, we write $\mathcal D(\pi)$ to denote $\mathcal D(\{\text{all vectors in $\pi$}\})$.
\end{defn}

Note $\psi$ in \cref{eqn:DMat} is distinct from the nonlinearities in the program $\pi$.
For example, if $\pi$ expresses the forward pass of a ReLU MLP, then $\pi$ has unbounded nonlinearity (ReLU) but $\psi$ in \cref{eqn:DMat} can be the step function, which is bounded.
We have kept $\psi$ bounded in \cref{eqn:DMat} for technical reasons (similar to the appearance of \emph{bounded} continuous functions in the definition of convergence in distribution).
But we believe \cref{thm:FreeIndependencePrinciple} holds when $\psi$ is more generally polynomially bounded.
This generalization would follow if \cref{thm:NetsorTMasterTheorem} holds for almost sure convergence of conditional means; see \cref{conj:NetsorTMeanConvergence}.

\paragraph*{Intuition and Discussion}
One can perhaps accept \netsort{} programs as a formalization of ``reasonable ways'' to compute vectors (and their diagonal matrices) from a set of random matrices.
Then FIP says that a random Gaussian matrix is asymptotically free from any diagonal matrix that ``depends on it in a reasonable way.''
This formalizes the intuition that singular vectors of Gaussian matrices are in general position to singular vectors of diagonal matrices, so one may expect these matrices to be asymptotically free.

We shall see next that \cref{thm:FreeIndependencePrinciple} allows us to easily compute the asymptotic Jacobian singular value distribution of a randomly initialized neural network. 

\paragraph{Extension to \netsortplus{} programs with variable dimensions}
\cref{thm:FreeIndependencePrinciple} holds as stated for \netsort{} programs with variable dimensions.
It also holds for \netsortplus{} programs with variable dimensions if nonlinearities are parameter-controlled and rank stability (\cref{assm:asRankStab}) is satisfied.

\end{document}